% =============================================================================
%  Notas de clase — Sesión 05: Bit Manipulation
%  Programación Avanzada — Universidad Panamericana
%  Compilar: pdflatex sesion05_bit_manipulation.tex
% =============================================================================
\documentclass[12pt, oneside]{article}
\usepackage{progavanzada}

\begin{document}

\portada{Sesión 05}{Bit Manipulation}{2026}

\tableofcontents
\newpage

% =============================================================================
\section{¿Por qué operar directamente sobre bits?}
% =============================================================================

Ya sabemos que toda la información en una computadora son bits.
Hasta ahora los hemos leído y convertido entre bases, pero sin tocarlos.
Esta sesión vamos a \textbf{manipularlos directamente} desde C++.

¿Para qué sirve eso?

\begin{itemize}
  \item \textbf{Velocidad}: las operaciones de bits son las más baratas que
        existen en un procesador — cuestan un solo ciclo de reloj, sin importar
        si el operando tiene 8 o 64 bits.
  \item \textbf{Compacidad}: puedes empacar 32 valores booleanos en un solo
        \texttt{int}, en lugar de usar 32 variables separadas.
  \item \textbf{Programación de sistemas}: drivers, protocolos de red,
        compresión, cifrado y programación competitiva hacen uso intensivo
        de estas operaciones.
\end{itemize}

Los operadores que vamos a estudiar son:

\begin{center}
\begin{tabular}{lll}
  \toprule
  \textbf{Operador} & \textbf{Símbolo} & \textbf{Qué hace} \\
  \midrule
  Desplazamiento izquierdo & \cpp{<<}                    & Mueve todos los bits a la izquierda \\
  Desplazamiento derecho   & \cpp{>>}                    & Mueve todos los bits a la derecha \\
  NOT bit a bit            & \texttt{\textasciitilde}    & Invierte cada bit \\
  AND bit a bit            & \texttt{\&}                 & 1 solo si ambos bits son 1 \\
  OR bit a bit             & \texttt{|}                  & 1 si al menos un bit es 1 \\
  XOR bit a bit            & \texttt{\textasciicircum}   & 1 si los bits son distintos \\
  \bottomrule
\end{tabular}
\end{center}

% =============================================================================
\section{Visualizar bits: \texttt{bitset<N>}}
% =============================================================================

Antes de operar, necesitamos \textbf{ver} los bits. C++ incluye \texttt{bitset<N>}
en la biblioteca \texttt{<bitset>}: convierte cualquier entero a su representación
en exactamente $N$ bits.

\begin{ejemplo}[Ver 30 en binario]
\begin{verbatim}
#include <bitset>
using bin8 = bitset<8>;

int x = 30;
cout << bin8(x);   // imprime: 00011110
\end{verbatim}
$30 = 16 + 8 + 4 + 2$, y los cuatro bits bajos encendidos confirman eso.
\end{ejemplo}

\texttt{bitset} es solo para \textbf{visualización}: no cambia el valor
almacenado ni el tipo de la variable.

\subsection{Literales en distintas bases}

C++14 permite escribir el mismo valor en decimal, binario (\cpp{0b}) o
hexadecimal (\cpp{0x}). Para el compilador son idénticos.

\begin{center}
\begin{tabular}{lll}
  \toprule
  \textbf{Forma} & \textbf{Ejemplo} & \textbf{Valor} \\
  \midrule
  Decimal        & \cpp{30}          & 30 \\
  Binario        & \cpp{0b00011110}  & 30 \\
  Hexadecimal    & \cpp{0x1E}        & 30 \\
  \bottomrule
\end{tabular}
\end{center}

Usar \cpp{0b} o \cpp{0x} no cambia nada en tiempo de ejecución — es solo
una comodidad para que el código sea más legible cuando trabajamos con máscaras
o valores relacionados con bits.

\begin{consejo}
  Un \cpp{char} también es un entero de 8 bits. Puedes inicializarlo con
  \cpp{char c = 0b01001101;} (que resulta ser la \texttt{'M'}) y usar
  \texttt{bitset} para ver sus bits igual que con un \texttt{int}.
\end{consejo}

% =============================================================================
\section{Patrones especiales}
% =============================================================================

Tres patrones aparecen constantemente en cualquier código que manipule bits.
Vale la pena memorizarlos.

\subsection{Todos los bits apagados: \texttt{x = 0}}

Obvio pero importante: el valor 0 tiene todos sus bits en 0.

\subsection{Todos los bits encendidos: \texttt{x = -1}}

En complemento a 2, $-1 = 2^n - 1$, que en binario es todo 1s sin importar
cuántos bits tenga el tipo. Es la forma idiomática en C++ de poner a 1 todos
los bits de una variable.

\begin{verbatim}
int x = -1;   // 11111111 11111111 11111111 11111111  (32 bits)
\end{verbatim}

\subsection{Un solo bit encendido: \texttt{1 << k}}

\cpp{1 << k} tiene exactamente el bit $k$ encendido y todos los demás apagados.
Vale $2^k$.

\begin{center}
\begin{tabular}{ccc}
  \toprule
  \textbf{Expresión} & \textbf{Valor} & \textbf{Binario (8 bits)} \\
  \midrule
  \cpp{1 << 0} & 1   & \texttt{00000001} \\
  \cpp{1 << 1} & 2   & \texttt{00000010} \\
  \cpp{1 << 2} & 4   & \texttt{00000100} \\
  \cpp{1 << 3} & 8   & \texttt{00001000} \\
  \cpp{1 << 4} & 16  & \texttt{00010000} \\
  \cpp{1 << 5} & 32  & \texttt{00100000} \\
  \cpp{1 << 6} & 64  & \texttt{01000000} \\
  \cpp{1 << 7} & 128 & \texttt{10000000} \\
  \bottomrule
\end{tabular}
\end{center}

Estas expresiones se llaman \textbf{máscaras de un bit} y son el ingrediente
básico de casi toda técnica de bit manipulation.

% =============================================================================
\section{Desplazamientos}
% =============================================================================

\subsection{Desplazamiento izquierdo \texttt{<<}}

\cpp{x << k} mueve todos los bits de \texttt{x} hacia la izquierda $k$
posiciones. Los bits que salen por la izquierda se \textbf{pierden}; los
huecos de la derecha se llenan con \textbf{0}.

\medskip
\textbf{Efecto aritmético:} \cpp{x << k} $= x \times 2^k$
(siempre que no haya desbordamiento).

\begin{ejemplo}[\texttt{6 << 2 = 24}]
\begin{center}
\begin{tabular}{ll}
  \texttt{6} (antes):  & \texttt{0 0 0 0 0 1 1 0} \\
  Desplazar 2:         & $\leftarrow\leftarrow$ \\
  \texttt{24} (después): & \texttt{0 0 0 1 1 0 0 0} \\
\end{tabular}
\end{center}
$6 \times 4 = 24$. \checkmark
\end{ejemplo}

\subsection{Desplazamiento derecho \texttt{>>}}

\cpp{x >> k} mueve todos los bits hacia la derecha $k$ posiciones.
Los bits que salen por la derecha se \textbf{pierden}; los huecos de la
izquierda se llenan con \textbf{0} (para enteros positivos).

\medskip
\textbf{Efecto aritmético:} \cpp{x >> k} $= \lfloor x / 2^k \rfloor$
(división entera, trunca hacia abajo).

\begin{ejemplo}[\texttt{100 >> 3 = 12}]
$100 / 8 = 12.5$ — el desplazamiento trunca y da 12.

En binario: \texttt{01100100} $\;\to\;$ \texttt{00001100} = 12. \checkmark
\end{ejemplo}

\begin{advertencia}
  El desplazamiento a la derecha de \textbf{negativos} es
  \textit{implementation-defined} en C++ hasta C++19.
  En la práctica, la mayoría de los compiladores extienden con el bit de
  signo (desplazamiento aritmético), pero no es obligatorio por el estándar.
  Para evitar sorpresas, trabaja con enteros sin signo (\texttt{unsigned int})
  cuando desplaces a la derecha.
\end{advertencia}

% =============================================================================
\section{Operadores lógicos bit a bit}
% =============================================================================

A diferencia de los operadores lógicos \texttt{\&\&} y \texttt{||} que
operan sobre valores completos, los operadores bit a bit actúan
\textbf{posición a posición} de forma independiente sobre cada par de bits.

\subsection{NOT bit a bit (\texttt{\textasciitilde})}

\texttt{\textasciitilde{}x} invierte cada bit: 0 $\to$ 1, 1 $\to$ 0.

\begin{center}
\begin{tabular}{cc}
  \toprule
  bit & \texttt{\textasciitilde}bit \\
  \midrule
  0 & 1 \\
  1 & 0 \\
  \bottomrule
\end{tabular}
\end{center}

\begin{recuerda}
  En complemento a 2, \texttt{\textasciitilde{}x} equivale siempre a $-(x+1)$:
  \[
    \mathtt{\sim 30} = -31, \qquad \mathtt{\sim 0} = -1, \qquad \mathtt{\sim(-1)} = 0
  \]
\end{recuerda}

\subsection{AND bit a bit (\texttt{\&})}

\texttt{a \& b} produce \textbf{1} solo si \textbf{ambos} bits
correspondientes son 1.

\begin{center}
\begin{tabular}{ccc}
  \toprule
  \texttt{a} & \texttt{b} & \texttt{a \& b} \\
  \midrule
  1 & 1 & 1 \\
  1 & 0 & 0 \\
  0 & 1 & 0 \\
  0 & 0 & 0 \\
  \bottomrule
\end{tabular}
\end{center}

\textbf{Uso típico:} \textit{apagar} bits específicos y \textit{consultar}
si un bit está encendido.

\begin{ejemplo}[\texttt{23 \& 35 = 3}]
\begin{center}
\begin{tabular}{ll}
  \texttt{23}:       & \texttt{0001 0111} \\
  \texttt{35}:       & \texttt{0010 0011} \\
  \texttt{23 \& 35}: & \texttt{0000 0011} = 3 \\
\end{tabular}
\end{center}
Solo sobreviven los bits encendidos en \textbf{los dos} operandos a la vez
(posiciones 0 y 1).
\end{ejemplo}

\subsection{OR bit a bit (\texttt{|})}

\texttt{a | b} produce \textbf{1} si \textbf{al menos uno} de los dos bits
es 1.

\begin{center}
\begin{tabular}{ccc}
  \toprule
  \texttt{a} & \texttt{b} & \texttt{a | b} \\
  \midrule
  1 & 1 & 1 \\
  1 & 0 & 1 \\
  0 & 1 & 1 \\
  0 & 0 & 0 \\
  \bottomrule
\end{tabular}
\end{center}

\textbf{Uso típico:} \textit{encender} bits específicos.

\begin{ejemplo}[\texttt{23 | 35 = 55}]
\begin{center}
\begin{tabular}{ll}
  \texttt{23}:      & \texttt{0001 0111} \\
  \texttt{35}:      & \texttt{0010 0011} \\
  \texttt{23 | 35}: & \texttt{0011 0111} = 55 \\
\end{tabular}
\end{center}
El resultado recoge todos los bits encendidos en cualquiera de los dos operandos.
\end{ejemplo}

\subsection{XOR bit a bit (\texttt{\textasciicircum})}

\texttt{a \textasciicircum{} b} produce \textbf{1} si los dos bits son
\textbf{distintos} (eXclusive OR — OR exclusivo).

\begin{center}
\begin{tabular}{ccc}
  \toprule
  \texttt{a} & \texttt{b} & \texttt{a \textasciicircum{} b} \\
  \midrule
  1 & 1 & 0 \\
  1 & 0 & 1 \\
  0 & 1 & 1 \\
  0 & 0 & 0 \\
  \bottomrule
\end{tabular}
\end{center}

\textbf{Uso típico:} \textit{invertir} bits específicos.

\begin{ejemplo}[\texttt{23 \textasciicircum{} 35 = 52}]
\begin{center}
\begin{tabular}{ll}
  \texttt{23}:                      & \texttt{0001 0111} \\
  \texttt{35}:                      & \texttt{0010 0011} \\
  \texttt{23 \textasciicircum{} 35}: & \texttt{0011 0100} = 52 \\
\end{tabular}
\end{center}
Los bits donde 23 y 35 \textbf{coinciden} se apagan; los que \textbf{difieren}
se encienden.
\end{ejemplo}

\begin{tecnico}[Propiedades del XOR]
  Dos identidades del XOR son fundamentales en algoritmos:
  \begin{itemize}
    \item \texttt{x \textasciicircum{} x = 0} siempre
          (cualquier valor XOR consigo mismo es cero).
    \item \texttt{x \textasciicircum{} 0 = x} siempre
          (XOR con cero deja el valor intacto).
  \end{itemize}
  Consecuencia: \texttt{(x \textasciicircum{} k) \textasciicircum{} k = x}.
  Aplicar XOR dos veces con el mismo valor regresa al original.
  Esta propiedad es la base de varios algoritmos de cifrado simple y de
  detección de errores.
\end{tecnico}

% =============================================================================
\section{Combinando operadores}
% =============================================================================

Los operadores de bits se combinan igual que la aritmética normal.
El orden de precedencia en C++ (de mayor a menor) es:

\begin{center}
\begin{tabular}{cl}
  \toprule
  \textbf{Precedencia} & \textbf{Operadores} \\
  \midrule
  Mayor & \texttt{\textasciitilde} (NOT) \\
        & \texttt{<<}, \texttt{>>} (desplazamientos) \\
        & \texttt{\&} (AND) \\
        & \texttt{\textasciicircum} (XOR) \\
  Menor & \texttt{|} (OR) \\
  \bottomrule
\end{tabular}
\end{center}

\begin{consejo}
  Cuando combines operadores, usa paréntesis generosamente.
  \texttt{a \& b | c} no siempre hace lo que parece a primera vista.
\end{consejo}

\begin{ejemplo}[Evaluar \texttt{(\textasciitilde(13 >> 2)) \& 39} paso a paso]
\begin{center}
\begin{tabular}{lll}
  \toprule
  \textbf{Paso} & \textbf{Expresión} & \textbf{Resultado} \\
  \midrule
  1 & \texttt{x = 13}                        & \texttt{0000 1101} \\
  2 & \texttt{x >> 2}                        & \texttt{0000 0011} = 3 \\
  3 & \texttt{\textasciitilde(x >> 2)}        & \texttt{...1111 1100} (invierte todo) \\
  4 & \texttt{y = 39}                        & \texttt{0010 0111} \\
  5 & \texttt{\textasciitilde(x>>2) \& y}    & \texttt{0010 0100} = \textbf{36} \\
  \bottomrule
\end{tabular}
\end{center}
El NOT produce un número con casi todos los bits encendidos; el AND final
con 39 filtra y deja pasar solo los bits que 39 tiene encendidos.
\end{ejemplo}

% =============================================================================
\section{Ejercicios}
% =============================================================================

\subsection{Desplazamientos}

\begin{enumerate}

  \item Sin usar calculadora, calcula el resultado decimal y la representación
        binaria de 8 bits de cada expresión.
        \begin{enumerate}
          \item \texttt{7 << 2}
          \item \texttt{100 >> 2}
          \item \texttt{1 << 5}
          \item \texttt{96 >> 3}
        \end{enumerate}

  \item ¿Cuántos desplazamientos a la izquierda necesitas para convertir
        el valor 3 en 48? Justifica.

  \item El valor de \texttt{x} es un número entre 1 y 255 (8 bits).
        Describe qué le ocurre al patrón de bits si aplicas \texttt{x >> 8}.

\end{enumerate}

\subsection{NOT, AND, OR y XOR}

\begin{enumerate}[resume]

  \item Calcula los siguientes resultados a mano. Muestra el proceso con
        la representación de 8 bits de cada operando.
        \begin{enumerate}
          \item \texttt{17 \& 12}
          \item \texttt{13 | 12}
          \item \texttt{17 \textasciicircum{} 12}
          \item \texttt{\textasciitilde{} 12}
                (expresa el resultado en decimal)
        \end{enumerate}

  \item El valor \texttt{x = 0b10110011}. Evalúa:
        \begin{enumerate}
          \item \texttt{x \& 0b00001111} \quad (¿qué hace esta operación?)
          \item \texttt{x | 0b00001111} \quad (¿qué hace esta operación?)
          \item \texttt{x \textasciicircum{} 0b11111111}
                \quad (¿qué hace esta operación?)
        \end{enumerate}

  \item Completa la tabla (todos los valores en 8 bits):
        \begin{center}
        \begin{tabular}{cccc}
          \toprule
          \texttt{a} & \texttt{b} & \texttt{a \& b} & \texttt{a | b} \\
          \midrule
          \texttt{11001100} & \texttt{10101010} & ? & ? \\
          \texttt{11110000} & \texttt{00001111} & ? & ? \\
          \texttt{10101010} & \texttt{10101010} & ? & ? \\
          \bottomrule
        \end{tabular}
        \end{center}

\end{enumerate}

\subsection{Combinando operadores}

\begin{enumerate}[resume]

  \item Evalúa las siguientes expresiones paso a paso, mostrando el resultado
        binario de 8 bits en cada etapa.
        \begin{enumerate}
          \item \texttt{(\textasciitilde{} 7) | (1 << 1) | (1 << 3) | (1 << 5)}
                seguido de \texttt{\& 29}
          \item \texttt{(\textasciitilde(13 >> 2)) \& 39}
        \end{enumerate}

  \item La expresión \texttt{x \textasciicircum{} y \textasciicircum{} y}
        siempre produce el mismo resultado sin importar el valor de \texttt{y}.
        ¿Cuál es ese resultado y por qué?

  \item Dado \texttt{n} un entero positivo, ¿qué representa la expresión
        \texttt{n \& (n - 1)} en términos de bits? Pruébalo con \texttt{n = 12}.

\end{enumerate}

% =============================================================================
\section{Resumen}
% =============================================================================

\begin{recuerda}
\begin{itemize}
  \item \textbf{\texttt{bitset<N>}}: visualiza $N$ bits de un entero.
        Solo para impresión, no cambia el valor.
  \item \textbf{Literales}: \cpp{0b...} (binario), \cpp{0x...} (hexadecimal);
        equivalentes en tiempo de ejecución.
  \item \textbf{Patrones clave}: \texttt{0} = todo apagado;
        \texttt{-1} = todo encendido;
        \cpp{1 << k} = solo el bit $k$ encendido ($= 2^k$).
  \item \textbf{Desplazamientos}:
        \cpp{x << k} $= x \times 2^k$;
        \quad \cpp{x >> k} $= \lfloor x / 2^k \rfloor$.
  \item \textbf{NOT} (\texttt{\textasciitilde{}x}):
        invierte todos los bits. Equivale a $-(x+1)$.
  \item \textbf{AND} (\texttt{\&}):
        1 si \textbf{ambos} bits son 1. Filtra/apaga bits.
  \item \textbf{OR} (\texttt{|}):
        1 si \textbf{al menos uno} es 1. Enciende bits.
  \item \textbf{XOR} (\texttt{\textasciicircum}):
        1 si los bits son \textbf{distintos}. Invierte bits.
        Propiedades: \texttt{x \textasciicircum{} x = 0};
        \quad \texttt{x \textasciicircum{} 0 = x}.
  \item \textbf{Precedencia} (mayor a menor):
        \texttt{\textasciitilde} $>$
        desplazamientos $>$
        \texttt{\&} $>$
        \texttt{\textasciicircum} $>$
        \texttt{|}. Usar paréntesis para claridad.
\end{itemize}
\end{recuerda}

\end{document}
